\chapter{مقدمه}

\section{انگیزه و اهمیت تحقیق}

دوره انتقال، شامل سه هفته پیش از زایش تا سه هفته پس از آن، به‌عنوان یکی از حساس‌ترین مراحل چرخه تولید در دامداری‌های صنعتی شناخته می‌شود، زیرا تغییرات فیزیولوژیکی شدید و وقوع بیماری‌های متابولیک و عفونی (مانند کتوز و جابه‌جایی شکمبه) می‌تواند سلامت دام و موفقیت دوره شیردهی را به خطر اندازد \cite{Drackley1999, LeBlanc2010, Salamone2022}. این چالش‌ها نه‌تنها بر تولید شیر تأثیر منفی می‌گذارند، بلکه با افزایش نرخ حذف دام‌ها و کاهش شاخص‌های اقتصادی، بار مالی قابل‌توجهی بر دامداری‌ها تحمیل می‌کنند. پیش‌بینی دقیق عملکرد شیر در اولین روز تست پس از زایش، ابزاری کلیدی برای ارزیابی موفقیت برنامه‌های مدیریت دوره انتقال و شناسایی زودهنگام مشکلات سلامتی است. با افزایش حجم داده‌های جمع‌آوری‌شده از دامداری ها، استفاده از الگوریتم‌های یادگیری ماشین برای تحلیل این داده‌ها و پیش‌بینی عملکرد شیر اهمیت روزافزونی یافته است \cite{Adriaens2021}.

برای ارزیابی کیفیت مدیریت در این دوره، شاخص دوره انتقال به‌عنوان ابزاری غیرمستقیم و مبتنی بر عملکرد تولیدی معرفی شده است که با مقایسه مقدار پیش‌بینی‌شده و واقعی اولین رکورد شیردهی، انحرافات ناشی از مدیریت نادرست را نشان می‌دهد \cite{Nordlund2006}. با این حال، روش‌های سنتی مانند تحلیل متابولیک خون، به دلیل هزینه بالا و خطاهای احتمالی، کارایی محدودی دارند \cite{DeKoster2019}. در مقابل، پیشرفت در جمع‌آوری داده‌های تاریخی و استفاده از تکنیک‌های یادگیری ماشین، امکان توسعه ابزارهای خودکار و دقیق برای پایش سلامت و پیش‌بینی تولید را فراهم کرده است \cite{GuevaraEscobar2022}. این ابزارها نه‌تنها به بهبود مدیریت گله کمک می‌کنند، بلکه می‌توانند در طراحی استراتژی‌های ژنتیکی، بهینه‌سازی تغذیه، و تصمیم‌گیری‌های کلان در صنعت دامداری نقش اساسی ایفا کنند \cite{Dematawewa2007}.

پژوهش حاضر با بهره‌گیری از داده‌های تاریخی و الگوریتم‌های نوین یادگیری ماشین، به دنبال ارائه مدلی برای پیش‌بینی اولین رکورد شیردهی گاوها است. این تلاش با توجه به تنوع عوامل مؤثر مانند شرایط محیطی و تولید دوره‌های پیشین، و نیاز به کنترل نویزهای غیرمدیریتی، از اهمیت ویژه‌ای برخوردار است و می‌تواند گامی مؤثر در جهت پایداری اقتصادی و سلامتی دام‌ها باشد.

\section{تعریف مسئله}

دوره انتقال، به دلیل پیچیدگی‌های زیستی و مدیریتی، یکی از چالش‌برانگیزترین مراحل در تولید شیر گاوهای شیری محسوب می‌شود. پیش‌بینی دقیق اولین رکورد شیردهی پس از زایمان، به‌عنوان ورودی کلیدی شاخص دوره انتقال، نیازمند مدلی است که بتواند داده‌های تاریخی، از جمله روزهای تست \LTRfootnote{Test Day Records - TDRs} و دوره‌های شیردهی پیشین، را با دقت تحلیل کند و تأثیر عوامل غیرمدیریتی (مانند فصل زایش، تعداد شکم زایش، و تولید دوره‌های قبلی) را از انحرافات ناشی از مدیریت تفکیک نماید \cite{Salamone2022, Dallago2019, Liseune2023}. روش‌های سنتی، مانند استفاده از منحنی‌های دوره شیردهی (لاکتاسیون) مانند مدل وود یا تحلیل متابولیک، به دلیل وابستگی به داده‌های محدود اولیه، هزینه‌های بالای آزمایشگاهی، و ناتوانی در در نظر گرفتن الگوهای فردی، اغلب ناکارآمد هستند \cite{GuevaraEscobar2022, Wood1967}.

در این پژوهش، مسئله اصلی توسعه الگویی داده‌محور است که با بهره‌گیری از داده‌های گسترده تاریخی روزهای تست و دوره‌های شیردهی پیشین، مقدار واقعی اولین رکورد شیردهی را پیش‌بینی کند. این مدل می‌تواند با استفاده از تکنیک‌های پیشرفته یادگیری ماشین، از جمله یادگیری عمیق، الگوهای پنهان در داده‌ها را استخراج کرده و به پیش‌بینی دقیق‌تر کمک کند \cite{Liseune2023}. تفاوت میان مقدار پیش‌بینی‌شده و واقعی، به‌عنوان شاخصی برای ارزیابی کیفیت مدیریت در سطح گله به کار می‌رود، به‌ویژه با میانگین‌گیری که نویزهای فردی را کاهش می‌دهد \cite{Nordlund2006}. با این حال، چالش‌هایی مانند کمبود داده در هفته‌های اولیه لاکتاسیون، تأثیر متغیرهای محیطی پیچیده، و نیاز به مدیریت نویزهای غیرمدیریتی، دقت مدل‌ها را تحت‌الشعاع قرار می‌دهد. بنابراین، استفاده از الگوریتم‌های انعطاف‌پذیر و قدرتمند، مانند شبکه‌های عصبی کانولوشنی، جنگل‌های تصادفی، و الگوریتم‌های مبتنی بر روش تقویت (مانند افزایش گرادیان افراطی و تقویت دسته‌ای)، که قادر به پردازش این پیچیدگی‌ها و بهبود پیش‌بینی‌ها باشند، ضروری به نظر می‌رسد \cite{Adriaens2021, Salamone2022}.

\section{اهداف تحقیق}

این پژوهش با هدف‌های زیر طراحی و اجرا شده است: ابتدا، انجام تحلیل و بررسی داده‌های یک دامداری با استفاده از ابزارهای آماری و ترسیم نمودارهای تحلیلی برای شناسایی الگوها و ویژگی‌های کلیدی داده‌ها، به‌منظور پشتیبانی از تصمیم‌گیری‌های مدیریتی انجام می‌شود. سپس، توسعه الگوریتم‌های مبتنی بر یادگیری ماشین، شامل رگرسیون ریج، جنگل‌های تصادفی، رگرسیون بردار پشتیبانی، افزایش گرادیان افراطی، تقویت دسته‌ای، و شبکه‌های عصبی، برای پیش‌بینی دقیق اولین رکورد شیردهی گاوهای شیری با بهره‌گیری از داده‌های دوره‌های شیردهی پیشین و روزهای آزمایش دنبال می‌شود. در نهایت، مقایسه عملکرد این مدل‌ها با استفاده از شش معیار ارزیابی، شامل ضریب تعیین، میانگین قدرمطلق خطا، ریشه میانگین مربع خطا، درصد خطای میانگین، میانگین قدرمطلق درصد خطای متقارن، و انحراف استاندارد نسبی، انجام می‌شود تا بهترین مدل برای پیش‌بینی انتخاب گردد. این فرآیند، علاوه بر بهبود دقت پیش‌بینی، به‌عنوان ابزاری برای ارزیابی غیرمستقیم کیفیت مدیریت دوره انتقال در سطح گله نیز به کار می‌رود.

\section{ساختار گزارش}

در فصل اول به بیان انگیزه، تعریف مسئله، اهداف، و ساختار کلی گزارش پرداخته می‌شود. در فصل دوم، مفاهیم پایه و کلیدی یادگیری ماشین، از جمله بیش‌برازش، کم‌برازش، و اعتبارسنجی متقاطع، و ...  به همراه توضیحات کامل مدل‌های به‌کاررفته در پژوهش، شامل رگرسیون ریج، جنگل‌های تصادفی، رگرسیون بردار پشتیبانی، افزایش گرادیان افراطی، تقویت دسته‌ای، شبکه‌های عصبی چندلایه، و شبکه‌های عصبی عمیق، ارائه می‌شود. همچنین، شش معیار ارزیابی شامل ضریب تعیین، میانگین قدرمطلق خطا، ریشه میانگین مربع خطا، درصد خطای میانگین، میانگین قدرمطلق درصد خطای متقارن، و انحراف استاندارد نسبی به همراه فرمول‌های مربوطه توضیح داده می‌شود. در ادامه، فصل سوم به بررسی مقالات مرتبط و رویکردهای موجود در پیش‌بینی تولید شیر اختصاص دارد. سپس، فصل چهارم ابتدا به تحلیل داده‌های یک دامداری با هدف آشنایی بهتر با مفهوم شاخص دوره انتقال می‌پردازد. پس از آن، مراحل آماده‌سازی داده‌ها، از جمله پاک‌سازی و سازمان‌دهی، و فرآیند آموزش و تنظیم پارامترهای مدل‌ها شرح داده می‌شود. در نهایت، فصل پنجم به مقایسه نتایج مدل‌ها با استفاده از معیارهای ارزیابی، تحلیل کمی و کیفی عملکرد آن‌ها، می‌پردازد.